% Plantilla del reporte de una prueba del enlace. reporte.py sustituye los
% campos @@...@@; no se compila directamente.
\documentclass[11pt]{article}
\usepackage[spanish,es-tabla,es-nodecimaldot,es-noshorthands]{babel}
\usepackage[utf8]{inputenc}
\usepackage[T1]{fontenc}
\usepackage{lmodern}
\usepackage[margin=2.5cm]{geometry}
\usepackage{amsmath}
\usepackage{graphicx}
\usepackage{xcolor}
\usepackage{booktabs}
\usepackage{tabularx}
\usepackage{float}
\usepackage{enumitem}
\usepackage[hidelinks]{hyperref}

\setlist{itemsep=2pt, topsep=4pt}
\pdfsuppresswarningpagegroup=1
\newcommand{\code}[1]{\texttt{#1}}

\title{Prueba del enlace PC--dsPIC\\\large @@TITULO@@}
\author{Haniel Ulises\\Departamento de Control Automático, Cinvestav}
\date{@@FECHA@@}

\begin{document}
\maketitle

\section{Configuración}

\begin{table}[H]
\centering
\begin{tabularx}{\linewidth}{@{}lX@{}}
\toprule
Montaje      & @@MODO@@ \\
Plataforma   & @@PLATAFORMA@@ \\
Reloj SPI    & @@RELOJ@@\,kHz, modo 0, \code{CS} activo en bajo \\
Ritmo        & @@RITMO@@ \\
Salida DAC   & @@SALIDA@@ \\
Comando      & \code{\small @@COMANDO@@} \\
@@FILA_VISOR@@
\bottomrule
\end{tabularx}
\end{table}

\section{Resultados}

Se enviaron @@TRAMAS@@ tramas en @@DURACION_TOTAL@@\,s (@@TASA@@ tramas/s).
Resultado: @@RESULTADO@@.

\begin{table}[H]
\centering
\caption{Errores por tipo. Un evento del dsPIC se reporta en la respuesta
siguiente a la trama que lo produjo.}
\begin{tabular}{@{}lrr@{}}
\toprule
Tipo & Tramas & \% \\
\midrule
@@FILAS_ERRORES@@
\bottomrule
\end{tabular}
\end{table}

\begin{table}[H]
\centering
\caption{Duración de \code{SPI\_ReadWrite} medida en la PC, en $\mu$s.}
\begin{tabular}{@{}rrrrrr@{}}
\toprule
Mínima & Media & Mediana & P99 & P99.9 & Máxima \\
\midrule
@@DT_MIN@@ & @@DT_MEDIA@@ & @@DT_MEDIANA@@ & @@DT_P99@@ & @@DT_P999@@ & @@DT_MAX@@ \\
\bottomrule
\end{tabular}
\end{table}

\begin{itemize}
@@OBSERVACIONES@@
\end{itemize}

@@NOTAS@@

\section{Gráficas}

@@FIGURAS@@

\end{document}
